\section{律师制度和律师职业道德}
\begin{enumerate}
  \item 跟公证员一起考。
  \item 行政管理。司法部，到县。
  \item 行业管理。全国，省一级。地级市可以有。
  \item 资格取得。过失犯罪可以。
        \begin{enumerate}
          \item 实习满 1 年。
          \item 要加入地方律协。只能一个律所执业。
          \item 可以兼职（人大代表）。省级发证。
        \end{enumerate}
  \item 律师事务所。
        \begin{enumerate}
          \item 3 年内没有受过停业处罚。
          \item 分所。3 年 20 人，省级审核。
          \item 无限连带。
          \item 关键信息，报批报备。特别关键，经过批准。一般重要，报告。
        \end{enumerate}
  \item 委托代理关系。基本同民法。
        \begin{enumerate}
          \item 利益冲突审查。直接利益冲突，必须解除；有风险的，主动提示风险 + 书面提示。
          \item 禁止虚假承诺。区分正常的诉讼风险。
          \item 保管委托人财产。
          \item 转委托。
          \item 可以终止。委托人隐瞒重大信息；没有履行 + 合理催告；委托人不切实际的目标。
        \end{enumerate}
  \item 律师执业/律所推广规范。
        \begin{enumerate}
          \item 回扣、夸大、暗示有关系。避免同行不正当竞争。
          \item 虚假宣传。不能出现\zd{专家}二字。
          \item 比较宣传。
        \end{enumerate}
  \item 收费。
        \begin{enumerate}
          \item 刑事案件、弱势群体、公共利益、国家赔偿。政府指导价。
          \item 律师费。
          \item 办案费。实报实销。
        \end{enumerate}
  \item 风险代理不能超过 18\%。政府指导价范围 + 人身性、伦理性。
  \item 同行之间\zd{是朋友}。
  \item 忠诚为民、法治正义、诚信敬业。\zd{只保障当事人合法利益}。
\end{enumerate}
